\chapter*{Abstrak}
\addcontentsline{toc}{chapter}{Abstrak}

\noindent\textit{Draf kolaborasi. Abstrak akan ditinjau setelah Bab 5--8 selesai.}\par\medskip

Tugas besar ini mengembangkan sistem pengenalan ucapan otomatis bahasa Sunda dengan menyetel ulang dua model wav2vec 2.0 XLSR-53 berbasis \textit{Connectionist Temporal Classification} (CTC). Korpus terdiri atas 2.185 ujaran dengan durasi sekitar 3 jam 32 menit dari 21 penutur. Himpunan latih, validasi, dan uji menggunakan penutur yang saling terpisah. Empat laju pembelajaran diuji, lalu bobot model dipilih berdasarkan \textit{Word Error Rate} (WER) pada data validasi. Setelah penyetelan, model Cahya memperoleh WER \pct{44,44} dan \textit{Character Error Rate} (CER) \pct{20,94} pada validasi, sedangkan model Indonesian NLP memperoleh WER \pct{41,59} dan CER \pct{20,41}. Pada 327 ujaran uji, penyetelan model Indonesian NLP menurunkan WER dari \pct{66,57} menjadi \pct{41,06} dan CER dari \pct{22,96} menjadi \pct{14,73}. Substitusi mendominasi kesalahan kata. Proporsi kata referensi yang salah lebih tinggi untuk kata yang tidak muncul dalam data latih, dan tingkat kesalahan berbeda antarpenutur. Evaluasi uji mencakup empat penutur dalam satu pembagian data.

\medskip
\noindent\textbf{Kata kunci} \enspace pengenalan ucapan otomatis, bahasa Sunda, wav2vec 2.0, CTC, WER, CER.
